\documentclass[../main.tex]{subfiles}

\begin{document}

\chapter{ Week 1+2 }

代靖涵 25120222201319

\begin{exercise}{}{}
A nonassociative algebra $(A, \circ)$ is called \dfntxt{left symmetric} if the associators are left symmetric in the sense that
\[\begin{aligned}
(x,y,z) = (y,x,z) \quad ((x,y,z) := (x \circ y) \circ z - x \circ (y \circ z))
\end{aligned}\]
for all $x,y,z \in A$. Prove that for every left symmetric algebra $(A, \circ)$, $(A, [\cdot, \cdot])$ is a Lie algebra, where
\[\begin{aligned}
[\cdot, \cdot] : A \times A \to A, \quad (x,y) \mapsto [x,y] := x \circ y - y \circ x
\end{aligned}\]
denotes the commutator on $A$.
\end{exercise}

\begin{proof}
    To show \(  \left[ \cdot ,\cdot  \right]   \) is bilinear, we  see  \[
    \begin{aligned}
    \left[ x+ y,z \right]&= \left( x+ y \right)\circ z- z\circ\left( x+ y \right)    \\ 
     &= \left( x\circ z - z\circ x \right) +\left(  y\circ z - z \circ y \right) \\ 
      &=  \left[ x,y \right]+ \left[ y,z \right]  
    \end{aligned}
    \]  Similarly \[
    \begin{aligned}
    \left[ x,y+ z \right]&= x\circ \left( y+ z \right)- \left( y+ z \right)\circ x\\ 
     &= \left( x\circ y- y\circ x \right)+ \left( x\circ z-z\circ x \right)\\ 
      &= \left[ x,y \right]+ \left[ y,z \right]        
    \end{aligned}
    \]

    Futhermore, it is obvious that \[
    \left[ x,x \right]= x\circ x-x\circ x= 0 
    \]
  Then we only need to show the Jacobi identity. 
 Since \(  \left( A,\circ  \right)   \) is left symmetric, we have 
    \[
    \begin{aligned}
    & \left[ x\left[ yz \right]  \right]+ \left[ y\left[ zx \right]  \right]+ \left[ z\left[ xy \right]  \right]   \\ 
     &= x\circ \left( y\circ z - z\circ y \right)- \left( y\circ z - z \circ y \right)\circ x\\ 
      &+  y\circ \left( z\circ x - x \circ z \right)-\left( z\circ x - x \circ z  \right)\circ y\\ 
       &+  z\circ \left( x\circ y-y\circ x \right)-\left( x\circ y-y\circ x \right)\circ z\\ 
        &= x\circ \left( y\circ z \right)- x\circ \left( z\circ y \right)-\left( y\circ z  \right)\circ x + \left( z\circ y \right)\circ x\\ 
         &+ y\circ \left( z\circ x \right)-y\circ \left( x\circ z \right)-\left( z\circ x  \right)\circ y             + \left( x\circ z \right)\circ y\\ 
          &+ z\circ \left( x\circ y \right) -z\circ \left( y\circ x \right)-\left( x\circ y \right)\circ z + \left( y\circ x \right)\circ z\\ 
           &=  \left( x\circ \left( y\circ z \right)-\left( x\circ y \right)\circ z   \right)- \left( x\circ \left( z\circ y \right)+ \left( x\circ z \right)\circ y   \right)\\ 
            &- \left( \left( y\circ z \right)\circ x  + y\circ \left( z\circ x \right) \right)+ \left( \left( z\circ y \right)\circ x-z\circ \left( y\circ x \right)   \right)\\ 
             &- \left( y\circ \left( x\circ z \right)- \left( y\circ x \right)\circ z   \right)-\left( \left( z\circ x \right)\circ y -z\circ \left( x\circ y \right)   \right)\\ 
              &= 0       
    \end{aligned}
    \] Thus \(  \left( A,\left[ \cdot ,\cdot  \right]  \right)    \) is a Lie algebra. 
\end{proof}

\begin{exercise}{}{}
Write $\operatorname{Der}(A)$ for the Lie algebra of derivations on $A$, where $A = \mathbb{C}[t, t^{-1}]$. Prove that $\operatorname{Der}(A)$ has a basis $\{L_n := t^{n+1} \frac{d}{dt} \mid n \in \mathbb{Z}\}$ and the Lie brackets among the basis elements are as follows:
\[\begin{aligned}
[L_n, L_m] = (m-n)L_{m+n} \quad (m,n \in \mathbb{Z}).
\end{aligned}\]
\end{exercise}

\begin{proof}
 
     Let \(  D\in \mathrm{Der}\left( A \right)   \) be any derivation on \(  A  \). Then \[
    \begin{aligned}
     D\left( 1  \right)&= D\left( t\cdot t^{-1}  \right) \\ 
      &=  D\left( t \right)\cdot t^{-1} + tD\left( t^{-1}  \right)   
    \end{aligned}  
     \]  The other hand \[
     D\left( 1 \right)= D\left( 1\cdot 1 \right)=  D\left( 1 \right)\cdot 1+ 1\cdot D\left( 1 \right)= 2D\left( 1 \right)   \implies D\left( 1 \right)= 0  
     \] Thus we have \[
     D\left( t \right)\cdot  t^{-1} + tD\left( t^{-1}  \right)= 0  \implies D\left( t^{-1}  \right)=  -t ^{-2}D\left( t \right) 
     \] From the linearity and the product rule, we observe that a dirivation \(  D \in \mathrm{Der}\left( A \right)   \) is totally determined by its value on  \(  t  \), that is  for any \(  D_1,D_2 \in \mathrm{Der}\left( A \right)   \) if \(  D_1\left( t \right)= D_2\left( t \right)    \), then \(  D_1= D_2  \).    It is straightforward to verify that  \(  L_{n}  \) is a derivation on \(  A  \) for each \(  n \in \mathbb{Z}   \). Note that \[
     L_{n}\left( t \right)=  t^{n+ 1} ,\quad  n\in \mathbb{Z} .
     \]   Since \(  D\left( t \right) \in A   \), we have  \[
     D\left( t \right)\in \operatorname{span}\left(\cdots ,L_{-2}\left( t \right),L_{-1}\left( t \right),L_0\left( t \right),L_1\left( t \right),L_2\left( t \right),\cdots       \right)  .
     \] Then we get \[
     D\in \operatorname{span}\left( \cdots ,L_{-2},L_{-1},L_{0},L_{1},L_{2},\cdots  \right) .
     \] Suppose we have a finite linear combination of the $L_n$ that equals the zero derivation:
$\sum_{n \in I} k_n L_n = 0,$
where $I \subset \mathbb{Z}$ is a finite set and $k_n \in \mathbb{C}$.
To check whether this is the zero derivation, we can test its action on the generator $t$.
$\left(\sum_{n \in I} k_n L_n\right)(t) = 0(t) = 0.$
This gives:
$\sum_{n \in I} k_n L_n(t) = \sum_{n \in I} k_n t^{n+1} = 0.$
The set of monomials $\{t^j\}_{j \in \mathbb{Z}}$ is a basis for $A$ as a $\mathbb{C}$-vector space, so they are linearly independent. The equation $\sum_{n \in I} k_n t^{n+1} = 0$ implies that all coefficients must be zero, i.e., $k_n = 0$ for all $n \in I$.
This proves that the set $\{L_n\}_{n \in \mathbb{Z}}$ is linearly independent. Combining above all , we know \(  \mathrm{Der}\left( A \right)   \) has a basis \(  \left\{ L_{n}:n\in \mathbb{Z}  \right\}  \).      

     To show \(  \left[ L_{n},L_{m} \right]= \left( m-n \right)L_{m+ n}    \), we only need to show that they agree on \(  t  \): \[
     \begin{aligned}
      \left[ L_{n},L_{m} \right]\left( t \right)&=  \left( L_{n} \left( L_{m}\left( t \right)  \right)-L_{m}\left( L_{n}\left( t \right)  \right)   \right)\\ 
       &=    \left( L_{n}\left( t^{m+ 1} \right)-L_{m}\left( t^{n+ 1} \right)   \right)\\ 
        &= \left( \left( m+ 1 \right)t^{m+ n+ 1}  -  \left( n+ 1 \right)t^{m+ n+ 1} \right)\\ 
         &= \left( m-n \right) t^{m+ n+ 1}\\ 
          &= \left( m-n \right)L_{m+ n}\left( t \right)     
     \end{aligned}
     \]   which completes the proof.

    
\end{proof}

\begin{exercise}{}{}
Let $\mathbb{F} = \mathbb{R}$ or $\mathbb{C}$, and let $V$ be an \dfntxt{n-dimensional} vector space over $\mathbb{F}$ equipped with a nondegenerate bilinear form
\[\begin{aligned}
B: V \times V \to \mathbb{F}, \quad (v,w) \mapsto B(v,w).
\end{aligned}\]
Prove that:
\begin{enumerate}
\item[(a)] If $\mathbb{F} = \mathbb{C}$ and $B$ is symmetric, then there exists a basis $\{v_1, v_2, \dots, v_n\}$ of $V$ such that $B(v_i, v_j) = \delta_{i,j}$ for all $1 \le i, j \le n$.
\item[(b)] If $\mathbb{F} = \mathbb{R}$ and $B$ is symmetric, then there exists a basis $\{v_1, v_2, \dots, v_n\}$ of $V$ and $p=0, 1, \dots, n$ such that $B(v_i, v_j) = \varepsilon_i \delta_{i,j}$ for all $1 \le i, j \le n$, where $\varepsilon_i = 1$ if $i \le p$ and $\varepsilon_i = -1$ if $i > p$.
\item[(c)] If $\mathbb{F} = \mathbb{R}$ or $\mathbb{C}$ and $B$ is skew-symmetric, then $n=2m$ is even and there exists a basis $\{v_1, v_2, \dots, v_n\}$ of $V$ such that the matrix
\[\begin{aligned}
(B(v_i, v_j))_{1 \le i, j \le n} = \begin{pmatrix} 0 & I_m \\ -I_m & 0 \end{pmatrix},
\end{aligned}\]
where $I_m$ denotes the identity matrix of rank $m$.
\end{enumerate}
\end{exercise}

\begin{proof}

     
     
     Lemma: If \(  V  \) is a vector space on the field \(  \mathbb{K}  \)  satisfying \(  \operatorname{char}\left( K \right)\neq 0   \), \(  B  \) is a nonzero symmetric bilinear form on \(  V  \) , then there exists \(  v \in V  \) such that \(  B\left( v,v \right)\neq 0   \).       

     Proof of Lemma:  Otherwise, we have for each \(  w,v\in V  \), \[
    \begin{aligned}
    0&= B\left( w+ v,w+ v \right)\\ 
     &= B\left( w,w \right)+ B\left( w,v \right)+ B\left( v,w \right)+ B\left( v,v \right)\\ 
      &= 2B\left( w,v \right)     
    \end{aligned}
    \]  which leads to a contradiction since \(  B  \) is nondegenerate. 


    \item (a) We induct on the dimension \(  n  \). If \(  n = 1  \),    there exists vector \(  v_0  \) such that \(  B\left( v_0,v_0 \right)\neq 0   \). We denote \(  B\left( v_0,v_0 \right)   \) by \(   \lambda  \in \mathbb{C}   \). \(   \lambda   \) has a square root in \(  C  \), denoted by \(  \sqrt{ \lambda }  \). Let \(  v_1= \frac{v_0 }{\sqrt{ \lambda } }   \). It constitute a basis of \(  V  \) such that  \[
    B\left( v_1,v_1 \right)= \frac{1 }{ \lambda  }B\left( v_0,v_0 \right)= 1   
    \]       Now we suppose that \(  \left( n-1 \right)   \)-dimensional space has property (a). We consider the \(  n  \)-dimensional vector space \(  V  \) on \(  \mathbb{C}   \), with a nondegenerate bilinear form \(  B  \).  By the Lemma, we can take a vector \(  v_0  \) such that \(  B\left( v_0,v_0 \right)\neq 0\) , and \(  v_1\coloneqq \frac{v_0 }{\sqrt{B\left( v_0,v_0 \right) } }   \) such that \(  B\left( v_1,v_1 \right)= 1   \).    Now we define \[
    W\coloneqq \left\{ w: B\left( v_1,w \right)= 0  \right\},
    \] which is a subspace of \(  V  \). We claim that \(  V= W\oplus \left<v_1 \right>  \), and thus \(  \operatorname{dim}W= n-1  \). 
    \begin{enumerate}
     \item \(  W\cap \left<v_1 \right>= 0  \):  Take any \(  w \in W\cap \left<v_1 \right>  \), then there is a complex valued \(  c  \) such that \(  w= cv_1  \). Futhermore,  \[
     0= B\left( w,v_1 \right)= cB\left( v_1,v_1 \right) = c\implies  c= 0 
     \]  There must be \(  w= 0  \).
     \item \(  W+ \left<v_1 \right>= V  \): Take any \(  v \in V  \) . Let \(  w=  v- B\left( v,v_1 \right)v_1   \), then \[
     B\left( w,v_1 \right)= B\left( v-B\left( v,v_1 \right)v_1,v_1  \right)=   B\left( v,v_1 \right)-B\left( v,v_1 \right)= 0  
     \]    Thus \(  w \in W  \), \(  v= w+ B\left( v,v_1 \right)v_1   \)   . Therefore \(  v \in W+ \left<v_1 \right>  \). 
    \end{enumerate}
    Now we prove that \(  B|_{W}: W\times W\to \mathbb{C}   \), \(  B|_{W}\left( w_1,w_2 \right)\coloneqq B\left( w_1,w_2 \right)    \) is a nondegenerate bilinear form on \(  W  \).
    \begin{enumerate}
     \item  Nondegenerate: For \(  w_0 \in W  \), if for any \(  w \in W  \), there is  \(  B|_{W}\left( w_0,w \right)= 0   \), then for any \(  v \in V  \), set \(  v=  k_1v_1+ k_2w_1  \) , \(  k_1,k_2\in \mathbb{C} , w_1\in W  \).  We have \[
     B\left( w_0,v \right) = k_1B\left( w_0,v_1 \right)+ k_2B\left( w_0,w_1 \right)= 0+ 0= 0  
     \]   Since \(  B  \) is nondegenerate, \(  w_0= 0  \). Therefore \(  B|_{W}  \) is nondegenerate.
     \item Symmetric and bilinear: Since \(  B  \) is symmetrc and bilinear, it is obvious that \(  B|_{W}  \) is as well.     
    \end{enumerate}
     From our inductive assumption, ther exists a basis \(  \left\{ v_2,\cdots ,v_{n} \right\}  \) for \(  W  \) such that \(  B\left( v_{i},v_{j} \right)=  \delta _{i,j}, 2\le i,j\le n   \). Then \(  \left\{  v_1,\cdots,v_n  \right\}   \) is a basis for \(  V  \) such that \(  B\left( v_{i},v_{j} \right)=  \delta _{i,j},1\le i,j\le n   \).      

     \item(b) We induct on the dimension \(  n  \). As in the (a), if \(  n = 1  \), there exists vector \(  v_0   \) such that \(  B\left( v_0,v_0 \right)\neq 0   \). We still denote \(  B\left( v_0,v_0 \right)   \) by \(   \lambda   \) , but here      \(   \lambda \in \mathbb{R}   \). We set \[
     v_1= \frac{v_0 }{\sqrt{\left|  \lambda  \right| } } 
     \]Then  \[
     B\left( v_1,v_1 \right)= B\left( \frac{v_0 }{\sqrt{\left|  \lambda  \right| } }, \frac{v_0 }{\sqrt{\left|  \lambda  \right| } }   \right)= \frac{ \lambda  }{\left|  \lambda  \right|  }= \begin{cases} 1,&  \lambda > 0\\ 
      -1,& \lambda < 0 \end{cases}    
     \]Now we suppose that any \(  \left( n-1 \right)   \)-dimensional space has property \(  \left( b \right)   \). We consider the \(  n  \)-dimensional vector space \(  V  \) on \(  \mathbb{R}   \). A smililar manner as above gives a vector \(u_1 \) such that \(  B\left( u_1,u_1 \right)= \pm 1 = : \varepsilon _{1}  \).    Still , define \[
     W\coloneqq \left\{ w:B\left( u_1 ,w \right)= 0  \right\}
     \]After a same process as in (a), we have \(  V= W\oplus \left<u_1 \right>  \), \(  \operatorname{dim}W= n-1  \), \(  B|_{W}  \) is a nondegenerate bilinear form on \(  W  \).   From our inductive assumption, there exists a basis \(  \left\{ u_2,\cdots ,u_{n} \right\}  \) for \(  W  \) such that \(  B\left( u_{i},u_{j} \right)=  \varepsilon _{i}  \delta _{i,j},  \varepsilon _{i}\in \left\{ -1,1 \right\}, 2\le i,j\le n   \)    . Then \(  \left\{ u_1,\cdots ,u_{n} \right\}  \) is a basis for \(  V  \) such that \(  B\left( u_{i},u_{j} \right)=  \varepsilon _{i} \delta _{i,j},  \varepsilon _{i}\in \left\{ -1,1 \right\}, 1\le i,j\le n   \).  Finally , we can relable \(  \left\{ u_1,\cdots ,u_{n} \right\}  \) to be \(  \left\{ v_1,\cdots v_{n} \right\}  \) such that there exists \(  p \in \left\{ 0,1,\cdots ,n \right\}  \) such that \(  B\left( v_{i},v_{j} \right)=  \varepsilon _{i} \delta _{i,j}   \) for all \(  1\le i,j\le n  \), where \(   \varepsilon _{i}= 1  \) if \(  i\le p  \) and \(   \varepsilon _{i}= -1  \) if \(  i> p  \).         

     (c) For each \(  v \in V  \), \(  B\left( v,v \right)= -B\left( v,v \right)\implies B\left( v,v \right)= 0     \) since \(  \operatorname{char}\left( \mathbb{F} \right)\neq 2   \) .   Let \(  V  \) be  vector space with arbitray positive dimension on \(  \mathbb{F}  \).  Claim that exists vector \(  v_0,w_0 \in V  \) such that \(  B\left( v_0,w_0 \right) \neq  0   \). Otherwise, for each \(  w,v \in V  \) ,  \[
     B\left( w,v \right)= B\left( v,w \right)= 0  
     \] , which is a contradiction to \(  B  \) is nondegenerate.  We set \[
     v_1= v_0,\quad w_1= \frac{w_0 }{B\left( v_0,w_0 \right)  } 
     \] Then \(  B\left( v_1,w_1 \right)= 1   \). Define \[
     W\coloneqq  \left\{ w: B\left( w, \operatorname{span}\left( v_1,w_1 \right)  \right)= 0  \right\}
     \]Then 
     \begin{enumerate}
          \item \(  W\cap \operatorname{span}\left( v_1,w_1 \right) = 0  \):  Take any \(  u \in W\cap \left<v_1,w_1 \right>  \), then there exists \(  m,n \in \mathbb{F}  \) such that  \(  u= mv_1+ nw_1  \). Then \[
          B\left( u,v_1 \right)= B\left( mv_1+ nw_1,v_1 \right)= nB\left( w_1,v_1 \right)= 0\implies n = 0   
          \]   \[
          B\left( u,w_1 \right)= B\left( mv_1+ nw_1,w_1 \right)= mB\left( v_1,w_1 \right)= 0\implies m= 0   
          \]Thus \(  u= 0  \). 
          \item \(  W+ \operatorname{span}\left( v_1,w_1 \right)= V   \)  : Take any \(  v \in V  \), let \(  w=  v- B\left(v,w_1 \right)v_1-B\left( v_1,v \right)w_1    \)  . Then \[
          B\left( w,v_1 \right)= B\left( v-B\left(v,w_1\right)v_1-B\left( v_1,v \right)w_1 ,v_1  \right)= B\left( v,v_1 \right)+ B\left( v_1,v \right)=    0  
          \] \[
          B\left(w,w_1 \right)= B\left( v-B\left( v,w_1 \right)v_1-B\left( v_1,v \right)w_1,w_1   \right)= B\left( v,w_1 \right)-   B\left(v,w_1 \right)= 0 
          \] Thus we have \(  V= W\oplus \operatorname{span}\left( v_1,w_1 \right)   \). \(  \operatorname{dim}W= \operatorname{dim}V-2  \).  

          Now we claim that the dimension of \(  V  \) must be even. Otherwise, once we get \(  V= W\oplus \operatorname{span}\left( v_1,w_1 \right)   \), we can do the same thing for \(  W  \) and \(  B|_{W}  \), and  get \(  W= W_1\oplus \operatorname{span}\left( v_2,w_2 \right)   \). But after finitely many decomposition, we get  \[
          V= \operatorname{span}\left( v_1,w_1 \right)\oplus \cdots \oplus \operatorname{span}\left( v_{m},w_{m} \right)\oplus W_{m-1}  
          \]    such that \(  \operatorname{dim}W_{m-1} = 1  \). Now there is no vector \(  v_{m},w_{m}   \) such that \(  B|_{W_{m-1}}\left( v_{m},w_{m} \right)= B\left( v_{m},w_{m} \right)> 0    \),   which is a contradiction. Thus there exists \(  m \in \mathbb{Z} _{> 0}   \) such that  \(  \operatorname{dim}V= 2m  \). We have shown above that there exists \(  v_1,\cdots ,v_{m}, w_1,\cdots ,w_{m} \in V  \) such that \[
          V= \operatorname{span}\left( v_1,w_1 \right)\oplus \cdots \oplus \operatorname{span}\left( v_{m},w_{m} \right)  
          \]   And \(  B\left( v_i,w_{i} \right)= 1,  1\le i\le m   \).  \(  B\left( v_{i},w_{j} \right)= 0, i\neq j   \).  We relabel \[
          \left\{v_1,\cdots ,v_{m} ,w_1,\cdots ,w_{m} \right\}
          \] to be \[
          \left\{ v_1,\cdots ,v_{m},v_{m+ 1},\cdots ,v_{n} \right\}
          \]Then we have \[
          B\left( v_{i},v_{i+ m} \right)= 1,\quad B\left( v_{i+ m},v_{i} \right)= -1  ,\quad i = 1,2,\cdots ,m
          \] and  \[
          B\left( v_{i},v_{j + m} \right) = 0,\quad  1\le i\neq j\le m
          \]Thus we have the  matrix
\[\begin{aligned}
(B(v_i, v_j))_{1 \le i, j \le n} = \begin{pmatrix} 0 & I_m \\ -I_m & 0 \end{pmatrix}.
\end{aligned}\]
     \end{enumerate}
     

\end{proof}
\end{document}